\section{Introduction}
\label{sec:intro}

\begin{guidance}[Introduction (10 marks), about half a page]
The marker wants: the context and goal of the project; your \textbf{selected
game} and \textbf{GUI library}, with the reason for each choice; your main
\textbf{design choices}; the \textbf{task allocation} between the two of you; and
the \textbf{structure} of the report. Keep it to about half a page (brief,
p.~6). Open with the problem, not with the technology. Expand every acronym at
first use, e.g.\ ``finite-state machine (FSM)''.
\end{guidance}

A serial game has two programs that can disagree about the state of the game.
If the PC is trusted, a faulty or hostile client can place two marks in one
turn, overwrite an occupied cell or keep playing after the game has ended, and
the board and the screen then show different games. The project brief~\cite{brief}
therefore makes the STM32 the referee. My design guarantees one property: a
request that the server rejects never changes the game state.

I chose Tic-Tac-Toe, the baseline game in the brief, so that my effort went
into the protocol, the robustness of the server and the automated tests rather
than into game rules. % CONFIRM: reword if you had a different reason.
The client uses Tkinter~\cite{tkinter} because it ships with the standard
Python installer and needs no extra installation, and pySerial~\cite{pyserial}
for the serial port. % CONFIRM: reword if you had a different reason.

Beyond the compulsory system I built three of the optional extensions that
the brief lists or recommends: a CRC-8 checksum on every frame, an AI opponent,
and move history with undo. Each one sits behind new commands, so the compulsory
protocol behaves exactly as before and the 23 core tests still apply.

Three design choices shape the system. First, the firmware splits reception in
two: a short interrupt handler stores each byte in a ring buffer, and the main
loop assembles lines, parses them and runs the state machine. Second, the
protocol is line-framed ASCII with no spaces, so a person can test the server
from a serial terminal. Third, one pytest harness drives the real board
without human interaction, and every robustness rule has a named test.

\textbf{Task allocation and plan.} The brief allows teams of two; I worked alone,
so I did all of the work: protocol, firmware, GUI, test harness, extensions and
this report. Table~\ref{tab:plan} lists the milestones in the order I worked.

\begin{table}[t]
\centering
\caption{Project plan: milestones and outputs.}
\label{tab:plan}
\footnotesize
\begin{tabular}{@{}L{0.3cm}L{4.9cm}L{2.2cm}@{}}
\toprule
 & Milestone and output & Date \\
\midrule
1 & Protocol and server FSM specified & 10/10/2026 \\
2 & Firmware: framing, state machine, validation & 16/10/2026 \\
3 & Tkinter GUI client & 16/10/2026 \\
4 & Test harness, 23 core tests; first board run & 16/10/2026 \\
5 & Receive-path fix after the failed tests & 24/10/2026 \\
6 & Extensions: checksum, AI, undo, 24 more tests & 24/10/2026 \\
7 & Final board runs, report, video, repository & 27/10/2026 \\
\bottomrule
\end{tabular}
\end{table}

\textbf{Structure.} Section~\ref{sec:protocol} defines the requirements, the
message protocol and the extensions. Section~\ref{sec:design} presents the server FSM, the GUI
flowchart and the circuit. Section~\ref{sec:impl} explains five code excerpts.
Section~\ref{sec:authority} argues that the server stays authoritative.
Section~\ref{sec:robust} covers framing and recovery from bad input.
Section~\ref{sec:valid} reports the tests on the board, and
Section~\ref{sec:conclusion} concludes.
